\documentclass[aspectratio=169, 10pt]{beamer}

\usetheme{default}
\useinnertheme{rounded}

% Couleurs
\definecolor{primary}{HTML}{0F172A}
\definecolor{secondary}{HTML}{1E40AF}
\definecolor{accent}{HTML}{6366F1}
\definecolor{success}{HTML}{059669}
\definecolor{warning}{HTML}{D97706}
\definecolor{danger}{HTML}{DC2626}
\definecolor{light}{HTML}{F1F5F9}
\definecolor{muted}{HTML}{64748B}

\setbeamercolor{structure}{fg=secondary}
\setbeamercolor{frametitle}{bg=primary, fg=white}
\setbeamercolor{block title}{bg=secondary, fg=white}
\setbeamercolor{block body}{bg=secondary!5, fg=primary}
\setbeamercolor{block title alerted}{bg=danger, fg=white}
\setbeamercolor{block body alerted}{bg=danger!5, fg=primary}
\setbeamercolor{block title example}{bg=success, fg=white}
\setbeamercolor{block body example}{bg=success!5, fg=primary}
\setbeamercolor{itemize item}{fg=accent}
\setbeamercolor{enumerate item}{fg=accent}
\setbeamercolor{normal text}{fg=primary}

\setbeamerfont{frametitle}{size=\large, series=\bfseries}
\setbeamertemplate{itemize item}{\small\raise0.5pt\hbox{\textbullet}}
\setbeamertemplate{navigation symbols}{}
\setbeamertemplate{footline}{
  \leavevmode%
  \hbox to \paperwidth{%
    \hspace*{2ex}\footnotesize\textcolor{muted}{Priora — Groupe Y}%
    \hfill%
    \footnotesize\textcolor{muted}{\insertframenumber{} / \inserttotalframenumber}%
    \hspace*{2ex}%
  }\vskip2pt%
}

\usepackage[utf8]{inputenc}
\usepackage[T1]{fontenc}
\usepackage[french]{babel}
\usepackage{booktabs}
\usepackage{tikz}
\usetikzlibrary{calc,positioning,shadows.blur}
\usepackage{colortbl}
\usepackage{tabularx}

\title{Priora}
\author[Groupe Y]{Adem DJEBBAR \and Nabil MOUSSAOUI \and Amir HADROUG \and Samy KHELLAF}

\begin{document}

% ===== TITRE =====
{
\setbeamercolor{background canvas}{bg=primary}
\begin{frame}[plain]
\vfill
\begin{center}
{\Huge\bfseries\textcolor{white}{Priora}}\par
\vspace{0.2cm}
{\normalsize\textcolor{accent!70!white}{Application de Gestion de Tâches avec Priorisation par IA}}\par
\vspace{0.7cm}
\textcolor{muted}{\rule{5cm}{0.3pt}}\par
\vspace{0.5cm}
{\small\textcolor{light}{Adem DJEBBAR \quad\textcolor{muted}{|}\quad Nabil MOUSSAOUI \quad\textcolor{muted}{|}\quad Amir HADROUG \quad\textcolor{muted}{|}\quad Samy KHELLAF}}\par
\vspace{0.5cm}
\begin{tikzpicture}
\node[rounded corners=3pt, fill=secondary, inner sep=5pt, font=\footnotesize\bfseries, text=white] {L3 Informatique — Génie Logiciel — Groupe Y};
\end{tikzpicture}\par
\vspace{0.2cm}
{\scriptsize\textcolor{muted}{Année universitaire 2025--2026}}
\end{center}
\vfill
\end{frame}
}

% ===== SOMMAIRE =====
\begin{frame}{Sommaire}
\tableofcontents
\end{frame}

% ================================================================
\section{Présentation du projet}
% ================================================================

\begin{frame}{Contexte et Problématique}
\begin{block}{Constat}
\small Les étudiants ont du mal à organiser leur travail au quotidien. Les applis de to-do classiques se contentent de lister les tâches sans aider à décider quoi faire en premier.
\end{block}
\vspace{0.3cm}
\textbf{\textcolor{secondary}{Notre idée :}} une application web qui utilise l'\textbf{IA} pour aider activement l'utilisateur :
\vspace{0.2cm}
\begin{itemize}
    \item Classer automatiquement les tâches par urgence et importance
    \item Conseils personnalisés basés sur les habitudes
    \item Planning journalier optimisé, adapté aux horaires
\end{itemize}
\end{frame}

\begin{frame}{Fonctionnalités — Utilisateurs et Tâches}
\begin{columns}[T]
\begin{column}{0.48\textwidth}
\begin{exampleblock}{Gestion des utilisateurs}
\small
\begin{itemize}
    \item Inscription avec mot de passe chiffré
    \item Connexion par token JWT
    \item Personnalisation des horaires et préférences
\end{itemize}
\end{exampleblock}
\end{column}
\begin{column}{0.48\textwidth}
\begin{block}{Gestion des tâches}
\small
\begin{itemize}
    \item Créer, modifier, supprimer des tâches
    \item Catégories : travail, études, personnel, autre
    \item Statuts : en attente, en cours, terminé
    \item Filtrer et trier
\end{itemize}
\end{block}
\end{column}
\end{columns}
\end{frame}

\begin{frame}{Fonctionnalités — IA et Notifications}
\begin{columns}[T]
\begin{column}{0.48\textwidth}
\begin{alertedblock}{Intelligence Artificielle}
\small
\begin{itemize}
    \item Priorisation selon la matrice d'Eisenhower
    \item Génération d'un planning journalier
    \item Analyse des habitudes et conseils
\end{itemize}
\end{alertedblock}
\end{column}
\begin{column}{0.48\textwidth}
\begin{block}{Notifications et alertes}
\small
\begin{itemize}
    \item Alerte quand des tâches sont en retard
    \item Résumé des tâches du jour
    \item Tableau de bord avec statistiques
\end{itemize}
\end{block}
\end{column}
\end{columns}
\end{frame}

% ================================================================
\section{Architecture technique}
% ================================================================

\begin{frame}{Technologies choisies}
\small
\begin{center}
\begin{tabular}{l l l}
\toprule
\textbf{Couche} & \textbf{Technologie} & \textbf{Pourquoi} \\
\midrule
Serveur & Python / FastAPI & Rapide, doc auto (Swagger) \\
\rowcolor{light}
Interface & React 18 / Vite & Composants, réactivité \\
Base de données & PostgreSQL & Robuste, relationnelle \\
\rowcolor{light}
ORM & SQLAlchemy + Alembic & Abstraction BDD + migrations \\
Auth & JWT + bcrypt & Sécurisé, stateless \\
\rowcolor{light}
IA & ChatGPT (OpenAI) & Langage naturel \\
Design & Tailwind CSS & Utilitaire, responsive \\
\bottomrule
\end{tabular}
\end{center}
\end{frame}

\begin{frame}{Architecture générale}
\begin{center}
\begin{tikzpicture}[
    node distance=2cm,
    box/.style={rectangle, thick, minimum height=1.1cm, minimum width=3.5cm, rounded corners=5pt, font=\footnotesize\bfseries, text width=3cm, align=center, blur shadow={shadow blur steps=5, shadow xshift=0.4mm, shadow yshift=-0.4mm}},
    arrow/.style={->, very thick, rounded corners}
]
    \node[box, draw=secondary, fill=secondary!10, text=secondary] (frontend) {Interface\\(React)};
    \node[box, draw=primary, fill=primary!10, text=primary, below=1.3cm of frontend] (api) {Serveur\\(FastAPI)};
    \node[box, draw=success, fill=success!10, text=success!80!black, below left=1cm and 1.5cm of api] (db) {Base de données\\(PostgreSQL)};
    \node[box, draw=accent, fill=accent!10, text=accent!80!black, below right=1cm and 1.5cm of api] (ai) {IA\\(ChatGPT)};

    \draw[arrow, color=secondary] (frontend) -- node[right, font=\scriptsize, text=muted] {HTTP / JSON} (api);
    \draw[arrow, color=success] (api) -- node[left, font=\scriptsize, text=muted] {lecture/écriture} (db);
    \draw[arrow, color=accent] (api) -- node[right, font=\scriptsize, text=muted] {prompt → JSON} (ai);
\end{tikzpicture}
\end{center}
\vspace{0.2cm}
\small L'utilisateur interagit avec React. Le serveur traite les requêtes, stocke dans PostgreSQL, et consulte ChatGPT pour les recommandations IA.
\end{frame}

\begin{frame}{Architecture en couches du backend}
\begin{center}
\begin{tikzpicture}[
    layer/.style={rectangle, thick, minimum height=0.9cm, minimum width=10cm, rounded corners=5pt, font=\footnotesize\bfseries, blur shadow={shadow blur steps=5, shadow xshift=0.4mm, shadow yshift=-0.4mm}},
    arrow/.style={->, thick, color=muted}
]
    \node[layer, draw=secondary, fill=secondary!10, text=secondary] (routers) at (0,2.5) {Couche Routes — reçoit et répond aux requêtes};
    \node[layer, draw=accent, fill=accent!10, text=accent!80!black] (services) at (0,1.25) {Couche Services — logique métier, IA, notifications};
    \node[layer, draw=success, fill=success!10, text=success!80!black] (models) at (0,0) {Couche Données — modèles SQLAlchemy, validation Pydantic};

    \draw[arrow] (routers) -- (services);
    \draw[arrow] (services) -- (models);
\end{tikzpicture}
\end{center}
\vspace{0.3cm}
\small Chaque couche a un seul rôle. On peut modifier la logique IA sans toucher aux routes, ou changer la BDD sans impacter le reste.
\end{frame}

% ================================================================
\section{Modélisation UML}
% ================================================================

\begin{frame}{Modèle de données — 3 entités}
\begin{columns}[T]
\begin{column}{0.32\textwidth}
\begin{block}{Utilisateur}
\scriptsize
Email, mot de passe chiffré, nom, prénom, horaires de travail. Entité centrale.
\end{block}
\end{column}
\begin{column}{0.32\textwidth}
\begin{block}{Tâche}
\scriptsize
Titre, catégorie, priorité, statut, date limite, durée estimée. Liée à un utilisateur.
\end{block}
\end{column}
\begin{column}{0.32\textwidth}
\begin{block}{Préférences}
\scriptsize
Catégories favorites, notifications, heure du résumé quotidien. 1 par utilisateur.
\end{block}
\end{column}
\end{columns}
\vspace{0.3cm}
\small\textbf{\textcolor{secondary}{Relations :}}
\begin{itemize}\small
    \item Utilisateur → Tâche : \texttt{1} vers \texttt{0..*} (possède plusieurs tâches)
    \item Utilisateur → Préférence : \texttt{1} vers \texttt{1} (un seul enregistrement)
\end{itemize}
\vspace{0.2cm}
\small\textbf{\textcolor{secondary}{3 énumérations :}} Catégorie (travail, études, personnel, autre), Priorité (basse à critique), Statut (en attente, en cours, terminé).
\end{frame}

\begin{frame}{Diagramme de cas d'utilisation}
\small
\textbf{\textcolor{secondary}{2 acteurs :}} Utilisateur et ChatGPT (service externe)
\vspace{0.2cm}

\textbf{\textcolor{secondary}{16 cas d'utilisation en 5 packages :}}
\vspace{0.1cm}

\footnotesize
\begin{tabular}{l l}
\toprule
\textbf{Package} & \textbf{Cas d'utilisation} \\
\midrule
\textcolor{secondary}{Authentification (4)} & S'inscrire, Se connecter, Profil, Se déconnecter \\
\rowcolor{light}
\textcolor{success}{Gestion des tâches (5)} & Créer, Modifier, Supprimer, Changer statut, Lister \\
\textcolor{accent}{Intelligence Artificielle (3)} & Prioriser, Générer planning, Analyser habitudes \\
\rowcolor{light}
\textcolor{warning}{Notifications (2)} & Tâches en retard, Résumé quotidien \\
\textcolor{danger}{Préférences (2)} & Configurer préférences, Modifier horaires \\
\bottomrule
\end{tabular}

\vspace{0.3cm}
\small\textbf{\textcolor{secondary}{Relations :}}
\begin{itemize}\small
    \item \texttt{<<include>>} (obligatoire) : Créer/Modifier/Prioriser incluent « Se connecter »
    \item \texttt{<<extend>>} (optionnel) : Changer statut étend Lister ; Horaires étend Préférences
\end{itemize}
\end{frame}

\begin{frame}{Diagramme de classes}
\small
\begin{columns}[T]
\begin{column}{0.48\textwidth}
\textbf{\textcolor{secondary}{3 entités :}}
\begin{itemize}\footnotesize
    \item \textbf{User} — email, mdp hashé, nom, horaires
    \item \textbf{Task} — titre, catégorie, priorité, statut, échéance
    \item \textbf{Preference} — catégories favorites, notifications
\end{itemize}
\vspace{0.2cm}
\textbf{\textcolor{secondary}{3 services :}}
\begin{itemize}\footnotesize
    \item \textbf{AIService} — priorisation, planning, analyse
    \item \textbf{AuthService} — hash mdp, JWT, auth
    \item \textbf{NotificationService} — retard, résumé jour
\end{itemize}
\end{column}
\begin{column}{0.48\textwidth}
\textbf{\textcolor{secondary}{Relations :}}
\begin{itemize}\footnotesize
    \item User \texttt{1 → 0..*} Task (possède)
    \item User \texttt{1 → 1} Preference (a)
    \item Services \texttt{..>} Entités (dépendance)
\end{itemize}
\vspace{0.2cm}
\textbf{\textcolor{secondary}{Visibilité :}}
\begin{itemize}\footnotesize
    \item \texttt{+} public (méthodes de service)
    \item \texttt{-} privé (attributs, méthodes internes)
\end{itemize}
\vspace{0.2cm}
\textbf{\textcolor{secondary}{3 enums :}}
\footnotesize Categorie, Priorite, Statut
\end{column}
\end{columns}
\end{frame}

% ================================================================
\section{Intelligence Artificielle}
% ================================================================

\begin{frame}{Le rôle de l'IA dans Priora}
\begin{block}{Comment ça va marcher}
\small Le serveur envoie les tâches à ChatGPT sous forme de prompt structuré. L'IA analyse et renvoie du JSON que le serveur transmet à l'interface.
\end{block}
\vspace{0.3cm}
\textbf{\textcolor{secondary}{Les 3 fonctions IA :}}
\begin{enumerate}\small
    \item \textbf{Priorisation} — Score d'urgence 0-10 par tâche (matrice d'Eisenhower) + justification
    \item \textbf{Planning} — Emploi du temps heure par heure adapté aux horaires et durées
    \item \textbf{Analyse} — Taux de complétion, catégorie la plus productive, conseils concrets
\end{enumerate}
\end{frame}

\begin{frame}{Exemple d'utilisation prévu}
\small Scénario concret :
\vspace{0.2cm}
\begin{enumerate}\small
    \item Un étudiant a 5 tâches dont 2 en retard et 1 deadline demain
    \item Il clique sur \textbf{« Prioriser mes tâches »}
    \item Le serveur construit un prompt et l'envoie à ChatGPT
    \item ChatGPT renvoie : deadline demain → score 9.2, retard → 8.5, etc.
    \item L'utilisateur voit le classement avec scores et justifications
\end{enumerate}
\vspace{0.3cm}
\small\textcolor{muted}{Gestion des erreurs prévue : si le service IA est indisponible, message clair à l'utilisateur.}
\end{frame}

% ================================================================
\section{Équipe et organisation}
% ================================================================

\begin{frame}{Les membres du groupe}
\small
\begin{center}
\begin{tabular}{l l l l}
\toprule
\textbf{Nom} & \textbf{N° étudiant} & \textbf{Rôle} & \textbf{Responsabilités} \\
\midrule
Adem DJEBBAR & 12304687 & \textcolor{secondary}{Backend Auth} & Modèles, BDD, JWT \\
\rowcolor{light}
Nabil MOUSSAOUI & 12203479 & \textcolor{secondary}{Backend Métier} & Tâches, intégration IA \\
Amir HADROUG & 12308871 & \textcolor{secondary}{Frontend} & Pages, composants, design \\
\rowcolor{light}
Samy KHELLAF & 12308991 & \textcolor{secondary}{Intégration} & Front/back, docs \\
\bottomrule
\end{tabular}
\end{center}
\vspace{0.3cm}
Tout le monde a participé aux décisions. Réunions \textbf{mardi et jeudi midi à la BU}. Communication via \textbf{Snapchat, WhatsApp et Git}.
\end{frame}

\begin{frame}{Planning du projet (4 semaines)}
\footnotesize
\begin{center}
\begin{tabular}{l l p{8cm}}
\toprule
\textbf{Semaine} & \textbf{Phase} & \textbf{Tâches prévues} \\
\midrule
\textcolor{secondary}{\textbf{Sem. 1}} & Conception & Cahier des charges, choix technos, diagrammes UML, dépôt Git \\
\rowcolor{light}
\textcolor{accent}{\textbf{Sem. 2}} & Fondations & BDD, modèles, auth JWT, projet React \\
\textcolor{success}{\textbf{Sem. 3}} & Fonctionnalités & CRUD tâches, intégration ChatGPT, design Tailwind \\
\rowcolor{light}
\textcolor{warning}{\textbf{Sem. 4}} & Finitions & Notifications, préférences, tests, documentation, soutenance \\
\bottomrule
\end{tabular}
\end{center}
\end{frame}

\begin{frame}{Défis anticipés}
\begin{columns}[T]
\begin{column}{0.48\textwidth}
\begin{alertedblock}{Intégration de l'IA}
\footnotesize S'assurer que ChatGPT renvoie du JSON exploitable. Prévoir un parsing robuste.
\end{alertedblock}
\vspace{0.15cm}
\begin{block}{Communication front/back}
\footnotesize React (port 5173) et FastAPI (port 8000) sur des ports différents — config CORS.
\end{block}
\end{column}
\begin{column}{0.48\textwidth}
\begin{alertedblock}{Sécurité}
\footnotesize Chiffrement des mots de passe (bcrypt) et gestion fiable des tokens JWT.
\end{alertedblock}
\vspace{0.15cm}
\begin{block}{Coordination équipe}
\footnotesize Définir les contrats d'interface (requêtes/réponses) pour avancer en parallèle.
\end{block}
\end{column}
\end{columns}
\end{frame}

% ================================================================
\section{Conclusion}
% ================================================================

\begin{frame}{Bilan et prochaines étapes}
\begin{columns}[T]
\begin{column}{0.48\textwidth}
\begin{exampleblock}{Accompli}
\small
\begin{itemize}
    \item Cahier des charges défini
    \item Technos et architecture choisies
    \item Modélisation UML complète
    \item Équipe organisée
    \item Dépôt Git en place
\end{itemize}
\end{exampleblock}
\end{column}
\begin{column}{0.48\textwidth}
\begin{block}{Prochaines étapes}
\small
\begin{itemize}
    \item Backend (auth, CRUD, service IA)
    \item Frontend (pages, composants, API)
    \item Intégration ChatGPT
    \item Tests et soutenance finale
\end{itemize}
\end{block}
\end{column}
\end{columns}
\end{frame}

% ===== FIN =====
{
\setbeamercolor{background canvas}{bg=primary}
\begin{frame}[plain]
\vfill
\begin{center}
{\Huge\bfseries\textcolor{white}{Merci !}}\par
\vspace{0.5cm}
{\large\textcolor{accent!70!white}{Des questions ?}}\par
\vspace{0.7cm}
\textcolor{muted}{\rule{3cm}{0.3pt}}\par
\vspace{0.3cm}
{\footnotesize\textcolor{muted}{Groupe Y — L3 Informatique — Génie Logiciel}}
\end{center}
\vfill
\end{frame}
}

\end{document}
